%% introduction --- the book's argument in one sitting: why chaos is worth
%% learning, what it is, and what the five parts do about it.
%%
%% Twin: frontmatter/en/introduction.tex.

\chapter*{\lblIntroduction}
\phantomsection
\addcontentsline{toc}{chapter}{\lblIntroduction}
\markboth{\lblIntroduction}{\lblIntroduction}

\noindent
Całkowite zaćmienie Słońca, które przeszło nad Hiszpanią 12 sierpnia 2026
roku, było w tablicach astronomów na długo przed urodzeniem kogokolwiek, kto
je oglądał. Jego pas, jego datę i minutę, w której się zaczęło, wyliczono z
ruchów Ziemi i Księżyca, a tablice zaćmień, liczone w ten sam sposób, sięgają
tysięcy lat naprzód.

Nikt nie potrafiłby powiedzieć ludziom czekającym w tym pasie, na miesiąc
przed zaćmieniem, czy niebo będzie na tyle czyste, żeby je zobaczyć.

Obie prognozy powstają z tego samego rodzaju matematyki. W obu przypadkach
komputer bierze obecny stan kawałka świata, stosuje reguły opisujące, jak się
on zmienia, i krok po kroku idzie naprzód w czasie. Reguły dla powietrza nie
są tajemnicą: to te same prawa ruchu i ciepła, które opisują Księżyc. A
jednak jedna prognoza jest dobra na tysiąclecia, a druga na mniej więcej
tydzień. Różnica nie polega na tym, że meteorolodzy są mniej staranni od
astronomów, ani na tym, że powietrze jest zbyt skomplikowane, by je
zapisać. Różnica ma nazwę, pomiar i dokładny opis matematyczny. Nazywa się
chaos i o nim jest ta książka.

\section*{Czym jest chaos}

W mowie potocznej chaos oznacza bałagan. W matematyce oznacza coś dużo
ciekawszego i niemal przeciwnego: układ, który podlega dokładnej regule, bez
żadnego przypadku, a mimo to zachowuje się w sposób, którego nie da się
przewidzieć na długo naprzód. Muszą być spełnione trzy warunki naraz:

\begin{itemize}[leftmargin=1.4em, itemsep=2pt]
  \item \textbf{reguła jest dokładna} \dash{} ten sam stan początkowy zawsze
        prowadzi do tej samej przyszłości, bez żadnej domieszki losowości;
  \item \textbf{małe różnice rosną} \dash{} dwa stany początkowe różniące się
        o mniej, niż cokolwiek da się zmierzyć, po pewnym czasie prowadzą do
        zupełnie różnych skutków;
  \item \textbf{ruch pozostaje w granicach} \dash{} nie ucieka do
        nieskończoności, więc rosnące różnice są wciąż zawracane do tego
        samego obszaru i nigdy się nie uspokajają.
\end{itemize}

Każdy z tych warunków jest prosty. Naukowców, którzy znaleźli je razem,
zaskoczyło to, jak często ta kombinacja występuje. Nie potrzebuje
skomplikowanej reguły: pierwszy układ chaotyczny w tej książce to jedna
linijka szkolnej algebry, która mieści się na kalkulatorze. Nie potrzebuje
losowości, tłumu ani tajemnicy. Pojawia się w kapiących kranach, wahadłach,
populacjach zwierząt, rytmie serca, orbitach planet i, najsłynniej, w
pogodzie.

\section*{Dlaczego warto poświęcić na to czas}

Możesz rozsądnie zapytać, po co ktoś, kto nie jest naukowcem, miałby się
tego uczyć. Oto powody, które przekonały mnie.

\emph{Żyjesz prognozami.} Aplikacje pogodowe, projekcje emerytalne, krzywe
epidemii, prognozy gospodarcze, szacunki korków i pewne siebie odpowiedzi
programów komputerowych to przewidywania tworzone przez modele. Teoria chaosu
mówi, którym z nich można w zasadzie ufać daleko naprzód, którym tylko przez
krótki czas i jak w każdym przypadku wygląda uczciwa postać prognozy.

\emph{Mierzy granicę, a ta granica jest dziwna.} W układzie chaotycznym błąd
stanu początkowego rośnie w każdym kroku mniej więcej o ten sam czynnik. To
znaczy, że tysiąckrotne ulepszenie pomiarów nie daje tysiąc razy dłuższej
prognozy; daje stały dodatkowy odcinek czasu, po którym mgła zapada tak samo
jak wcześniej. Kiedy raz zobaczysz to przeliczone na liczbach \dash{} a
zobaczysz, w części~II \dash{} już nigdy nie usłyszysz zdania
\enquote{potrzebujemy tylko lepszych danych} w ten sam sposób.

\emph{Wyjaśnia, dlaczego dobre prognozy podają szanse.} Jeśli dalekiej
przyszłości nie da się policzyć jako jednej odpowiedzi, to da się policzyć
rozrzut odpowiedzi i to, jak prawdopodobna jest każda z nich.
\enquote{70 procent szans na deszcz} to nie asekuracja prognosty. To
poprawna odpowiedź na pytanie, na które chaos nie pozwala odpowiedzieć
\enquote{tak} albo \enquote{nie}.

\emph{Rozdziela rzeczy, które zwykle wrzuca się do jednego worka.}
\enquote{Nie da się tego przewidzieć} może znaczyć, że nikt nie zna reguł,
że reguły zawierają przypadek, albo że reguły są znane dokładnie, a
przyszłość i tak jest poza zasięgiem. To różne sytuacje, wymagające różnych
reakcji, a chaos to ta trzecia.

\emph{Jest piękny i uniwersalny.} Kształty, które rysuje chaos, to fraktale,
struktury pełne szczegółów przy każdym powiększeniu, a niektóre jego liczby
okazują się takie same w układach, które nie mają ze sobą nic wspólnego.
Ciecz, obwód elektroniczny i linijka algebry mogą dzielić stałą, której nikt
tam nie włożył.

\emph{I jest użyteczny.} Ta sama wrażliwość, która czyni chaos
nieprzewidywalnym, pozwala nim sterować bardzo małymi pchnięciami, a układy
chaotyczne można zmusić, żeby szły ze sobą w parze. Część~V pokazuje jedno
i drugie.

\section*{Czym ta książka jest, a czym nie jest}

Nie jest podręcznikiem i nie zrobi z ciebie matematyka. Nie jest też książką
z obrazkami, z której usunięto matematykę. Każdy pomysł w niej to prawdziwy
pomysł, rozwinięty od arytmetyki w górę, a każde twierdzenie, które może
sprawdzić komputer, zostało przez komputer sprawdzone; programy są
wydrukowane w książce i możesz je uruchomić sam. Tam, gdzie książka coś
stwierdza bez dowodu, mówi to otwarcie i mówi, gdzie ten dowód jest.

Zaczyna od zera. Część~I (rozdziały od~\ref{ch:models}
do~\ref{ch:motion}) pokazuje, jak matematyka w ogóle opisuje świat: czym
jest model, jak regułę stosuje się krok po kroku w czasie, dlaczego sprzężenie
zwrotne nie pozwala niczemu rosnąć w nieskończoność i jak opisuje się ruch.
Część~II (rozdziały od~\ref{ch:logistic} do~\ref{ch:bifurcations}) znajduje
chaos w jednej linijce algebry, mierzy, jak szybko rosną małe błędy, zamienia
to w jedną liczbę i horyzont prognozy, i idzie drogą od porządku do chaosu.
Część~III (rozdziały od~\ref{ch:lorenz} do~\ref{ch:mandelbrot}) przygląda się
kształtom, jakie tworzy chaos: atraktorowi Lorenza, atraktorom dziwnym,
fraktalom i zbiorowi Mandelbrota. Część~IV (rozdziały od~\ref{ch:mechanics}
do~\ref{ch:computers}) wyprowadza go w świat: wahadła i planety, prognoza
pogody, układy żywe i rynki, i sam komputer. Część~V
(rozdziały~\ref{ch:taming} i~\ref{ch:why}) pokazuje, jak chaos można
kontrolować, i kończy się przewodnikiem terenowym do oceniania każdej
prognozy, jaką ktoś ci przedstawi.

\section*{Pytanie}

W 1814 roku matematyk Pierre-Simon Laplace wyobraził sobie inteligencję,
która zna każdą siłę w przyrodzie i dokładny stan wszystkiego w jednej
chwili. Dla takiej inteligencji, pisał, nic nie byłoby niepewne: przyszłość
byłaby równie jasna jak przeszłość. To najbardziej pewne siebie zdanie, jakie
kiedykolwiek napisano o przewidywaniu, i przez sto lat wyglądało na naturalny
wniosek z fizyki.

Ta książka pyta, czy miał rację, i odpowiada liczbami, które możesz policzyć
sam. Rozdział~\ref{ch:models} zaczyna się od kamienia i kubka herbaty.
